% ---------------------------------------------------------------------------
\section{Introduction}
% ---------------------------------------------------------------------------
Autonomous, closed-loop DDoS mitigation is a basic capability for zero-touch
6G service management~\cite{benzaid2020zsm,liyanage2023oran}. It is also a
setting where purely neural control is hard to trust: exploratory actions can
cause outages, learned policies are opaque, and policies trained against
adaptive opponents may not converge~\cite{guo2020xai6g,senevirathna2025xaib5gsec}.
Sentinel~\cite{alfatemi2026sentinel} is a recent and carefully evaluated
answer to these problems. It wraps a PPO defender~\cite{schulman2017ppo} in a
symbolic safety shield~\cite{alshiekh2018shielding}, trains it in a co-evolutionary
wargame against a PPO attacker with a Hall-of-Fame archive~\cite{rosin1997newmethods},
and distils the result into depth-4 decision trees. Its evaluation is also
unusually candid: all ten seeds show \emph{late-stage co-evolutionary
collapse}, so the reported numbers depend on a benchmark checkpoint selected
after training. Outside the held-out ICMP scenario, the shielded policy is
statistically indistinguishable from its unshielded PPO ablation, and strict SLA
compliance is never reached under attack.

This paper asks whether the \emph{co-evolutionary} part of the design is needed
at all. We propose \textbf{Neuro-Symbolic Continuous Solution Discovery
(NS-CSD)}, which keeps Sentinel's POMDP, action space and shield but replaces
its training paradigm. The equilibrium-seeking attacker--defender game becomes
an open-ended discovery process that accumulates solutions, both neural and
symbolic, and never discards the best one found so far. The defender is trained with Group
Relative Policy Optimization (GRPO)~\cite{shao2024deepseekmath}. GRPO is
critic-free: advantages come from comparing a group of rollouts that start
from the same state. In a simulator the group can be \emph{forked} from one
state and made to replay identical traffic (common random numbers,
CRN~\cite{ng2000pegasus}), so the group comparison isolates the effect of the
defender's actions. Scenario difficulty comes from a regret-driven curriculum
grounded in an archive of known solutions~\cite{jiang2021plr,parkerholder2022accel},
not from a learned adversary. Two further changes make the method
neuro-symbolic in both directions: symbolic programs are archive elites that
teach the neural policy, and the parameters of the symbolic shield are
themselves discovered under a safety certificate.

We re-implemented Sentinel's released simulator and shield as a vectorised,
test-checked port and re-evaluated Sentinel's own paper checkpoints with it.
This reproduced the published results closely and exposed three properties of
the testbed that bound what any defender can achieve (Section~\ref{sec:diagnostics}).
First, link overload is computed on the pre-filter load, which caps service
quality at 0.556 under a full-intensity flood whatever the defender does.
Second, the fixed shield's PanicGuard clamps mitigation whenever utilisation is
below 0.6, so low-rate floods pass almost untouched. Third, the shaped reward
used for training differs from the composite score used for model selection.
These observations shaped NS-CSD's design and the way we report results.

\textbf{Contributions.}
\begin{itemize}
\item \emph{A co-evolution-free training paradigm} for neuro-symbolic
DDoS defence. It combines CRN-forked, critic-free GRPO; a quality-diversity
Solution Archive in which neural snapshots, crystallised tree programs and seed
programs compete per scenario niche with monotone elitism; archive-regret
scenario discovery; and archive-guided improvement (elite injection, champion
KL anchor and roll-back). Model selection happens inside training, not after it.
\item \emph{Certificate-gated shield discovery.} The symbolic layer is
re-parameterised during training using shield probes inside GRPO groups and a
group-relative evolution-strategy direction. A change is adopted only if it
causes no safety regression relative to the operator-approved shield on
attack-free and flash-crowd traffic.
\item \emph{A reproducible testbed audit.} We provide a vectorised port of
Sentinel's simulator and shield that is equivalence-tested against the
reference code and reproduces the published Table~I from the released
checkpoints, together with diagnostics that separate simulator ceilings from
learning failures.
\item \emph{A three-protocol evaluation} against Sentinel's released 10-seed
checkpoints and an equal-budget PPO control. Protocol A replays Sentinel's own
trained attackers, which is a transfer setting for NS-CSD. Protocol B is an
attacker-agnostic battery with four held-out stress tests. Protocol C is an
adaptive cross-entropy red team over the full scenario space, ICMP included.
The evaluation includes ablations of every component, paired statistics with
Holm correction, and an analysis of crystallised rule programs.
\end{itemize}

% ---------------------------------------------------------------------------
\section{Related Work}
% ---------------------------------------------------------------------------
\textbf{Learning-based DDoS defence.} Deep detectors such as
LUCID~\cite{doriguzzi2020lucid} classify traffic but do not decide how to
mitigate it. DRL mitigation has used DDPG flow control in
SDN~\cite{liu2018drlddos}, multi-agent router throttling~\cite{malialis2015distributed},
per-host drop control~\cite{simpson2020perhost} and autonomous SDN mitigation
frameworks~\cite{akbari2020atmos}; see~\cite{nguyen2023drlcyber} for a survey.
Adversarial RL attackers are also used to stress detectors~\cite{shao2025adados}.
Sentinel~\cite{alfatemi2026sentinel} adds runtime shielding, rule
crystallisation and co-evolution; NS-CSD is built directly on its formulation.

\textbf{Critic-free policy optimisation.} GRPO~\cite{shao2024deepseekmath}
replaces PPO's learned value baseline~\cite{schulman2017ppo,schulman2016gae}
with a group mean, much like leave-one-out REINFORCE~\cite{kool2019buy4,ahmadian2024back}.
It underlies large-scale RL with verifiable rewards~\cite{guo2025deepseekr1},
and its biases and stabilisation have been studied~\cite{yu2025dapo,liu2025understanding,zheng2025gspo}.
LUFFY mixes off-policy traces into GRPO groups~\cite{yan2025luffy}. Other work
constrains GRPO with Lagrangian costs~\cite{girgis2026constrainedgrpo} or adds
diversity to its reward~\cite{li2025darling}. Outside language models, GRPO
has been studied for classic control and robotics~\cite{oliveira2025critics,khanda2025grpocontinuous,wang2026gr2po}.
It has also been applied to network security, as an LLM-based security-automation loop in
zero-touch networks~\cite{cao2025secloop}. We therefore do not claim the first
use of GRPO in either setting. To our knowledge, NS-CSD is the first use of
GRPO for a numeric DDoS mitigation controller. It is also the first to fork
CRN groups from simulator states and to couple GRPO with an archive of mixed
neural and symbolic solutions and with UED-style scenario generation.

\textbf{Curricula and open-endedness.} Unsupervised environment design
generates training levels by regret: PAIRED~\cite{dennis2020paired},
PLR~\cite{jiang2021plr,jiang2021robustplr}, ACCEL~\cite{parkerholder2022accel},
with a critique of regret proxies in~\cite{rutherford2024noregrets}.
POET co-evolves environments and solutions~\cite{wang2019poet}, and
curricula have been paired with GRPO for LLM prompts~\cite{chen2025sec}. Quality
diversity keeps a map of elites~\cite{mouret2015mapelites,pugh2016qd}, which
can be improved with policy gradients~\cite{nilsson2021pgame}. Program-discovery
systems such as FunSearch~\cite{romeraparedes2024funsearch} and
AlphaEvolve~\cite{novikov2025alphaevolve} keep monotone databases of
evaluated programs. These systems are the closest analogue of our continuous
solution discovery, which fits the open-endedness framing
of~\cite{hughes2024openendedness}.

\textbf{Co-evolution pathologies.} Competitive co-evolution is known to cycle,
disengage and forget~\cite{ficici1998challenges,watson2001minimal,balduzzi2019openended}.
Hall-of-Fame archives~\cite{rosin1997newmethods}, fictitious self-play~\cite{heinrich2015fsp}
and PSRO~\cite{lanctot2017psro} reduce these problems but keep a learned
opponent; the same holds for robust adversarial RL~\cite{pinto2017rarl}.
NS-CSD removes the opponent entirely.

\textbf{Safe and interpretable RL.} Shields enforce specifications at
runtime~\cite{alshiekh2018shielding,konighofer2021online}. Constrained RL moves
safety into the objective instead~\cite{achiam2017cpo,garcia2015saferl}.
VIPER~\cite{bastani2018viper} and programmatic RL~\cite{verma2018pirl} extract
verifiable programs from policies; see~\cite{acharya2023neurosymbolic} for a
survey of neuro-symbolic RL. XAI is widely argued to be necessary for 6G
automation~\cite{guo2020xai6g,wang2021xai6g}. NS-CSD differs from prior shielded
RL in treating the shield as a \emph{discoverable} artefact, while keeping a
runtime guarantee envelope defined by the operator.
